% !TeX root=../main.tex
\chapter{روش تحقیق}
\label{chap:problem}
%\thispagestyle{empty}

% ============================================================
%  مقدمه
% ============================================================
\section{مقدمه}

این فصل مشخصات کامل ریاضی مسئله‌ی تخصیص بودجه در سطح سبد پروژه‌ها و
پیاده‌سازی آن به‌صورت یک محیط یادگیری تقویتی سازگار با \lr{Gymnasium}
را ارائه می‌کند. بخش~\ref{subsec:problem_statement} دامنه و مفروضات
ساختاری مسئله را بیان می‌کند. بخش~\ref{subsec:environment} عناصر اولیه‌ی
محیط را تعریف می‌کند: دینامیک پیشرفت، سازوکارهای جریان نقدی
\lr{FIDIC}، سیگنال‌های عملکرد \lr{EVM}، سازوکار فسخ و موجودی خزانه.
بخش~\ref{subsec:parameters} همه‌ی پارامترها را یک‌جا گردآوری می‌کند.
بخش~\ref{subsec:theoretical_ub} کران بالای اطلاعات‌کامل را به‌دست
می‌آورد و نشان می‌دهد چرا حل دقیق مسئله‌ی تصادفی در مقیاس سبد از نظر
محاسباتی شدنی نیست. بخش~\ref{subsec:mdp} صورت‌بندی فرایند تصمیم مارکوف
(\lr{MDP}) را ارائه می‌دهد: فضای حالت، فضای عمل، تصویرسازی عمل، هسته‌ی
گذار، تابع پاداش و بردار مشاهده‌ای که رابط محیط با هر سیاستی است.

% ============================================================
%  محتوا
% ============================================================
\section{تدوین مسئله‌ی تخصیص بودجه در سبد پروژه و فرمول‌بندی آن در قالب فرایند تصمیم مارکوف}

\subsection{علت انتخاب روش}

مسئله‌ی مورد مطالعه، یعنی تخصیص پویای بودجه میان چند پروژه‌ی فعال تحت
سازوکارهای قراردادی \lr{FIDIC} و عدم‌قطعیت در کارایی، ذاتاً یک مسئله‌ی
تصمیم‌گیری ترتیبی تحت عدم‌قطعیت است. راه طبیعی تحقیق در عملیات برای
چنین مسئله‌ای، برنامه‌ریزی تصادفی چندمرحله‌ای است~\citep{Birge2011}.
اما همان‌طور که در بخش~\ref{subsubsec:intractability} به‌تفصیل نشان
داده می‌شود، اندازه‌ی درخت سناریوها در این روش با افزایش تعداد پروژه‌ها
و دوره‌های زمانی به‌صورت نمایی رشد می‌کند (رابطه‌ی~\eqref{eq:tree_size}).
برای مثال، در یک سبد با ۱۰ پروژه تعداد برگ‌های درخت سناریو به حدود
$10^{167}$ می‌رسد، عددی که از تعداد اتم‌های جهان قابل مشاهده بسیار بیشتر
است. این رشد ترکیباتی از خود صورت‌بندی برنامه‌ریزی تصادفی چندمرحله‌ای
می‌آید و با عوض‌کردن حل‌کننده از بین نمی‌رود.

به همین دلیل، روش تحقیق این پژوهش \textbf{تحلیلی و مدل‌سازی} است. مسئله
به‌صورت یک فرایند تصمیم مارکوف (\lr{MDP}) صورت‌بندی می‌شود. در این
صورت‌بندی به‌جای شمارش صریح سناریوها، سیاست تصمیم‌گیری با تجربه‌ی
مسیرهای نمونه‌برداری‌شده بهینه می‌شود و هزینه‌ی محاسباتی هر گام به
اندازه‌ی فضای سناریو بستگی ندارد. این انتخاب دو هدف را هم‌زمان برآورده
می‌کند. هدف نخست حفظ دقت ریاضی صورت‌بندی است، با تعریف رسمی \lr{MDP} که
فضای حالت، فضای عمل، هسته‌ی گذار و تابع پاداش آن کاملاً مشخص است. هدف
دوم فراهم‌کردن بستری است که با کتابخانه‌های استاندارد یادگیری تقویتی
سازگار باشد~\citep{Towers2024}، تا نتایج این پژوهش در کارهای آینده
مستقیماً برای آموزش عامل‌های یادگیری‌شده به‌کار برود. جزئیات این
صورت‌بندی و استدلال مربوط به شدنی‌نبودن روش دقیق در زیربخش بعد آمده است.

\subsection{تشریح کامل روش تحقیق}

% ------------------------------------------------------------
\subsubsection{بیان مسئله و دامنه}
\label{subsec:problem_statement}
% ------------------------------------------------------------

مسئله در سطح راهبری راهبردی تعریف می‌شود. سبدی از $N$ پروژه‌ی
قراردادشده ثابت است: تعهدها امضا شده‌اند، تاریخ‌های شروع زمان‌بندی
شده‌اند و قراردادها طوری پارامتری شده‌اند که با نسخه‌های کتاب قرمز،
زرد و نقره‌ای \lr{FIDIC} سازگار باشند~\citep{FIDIC2017}. مدیر سبد در هر
دوره‌ی گسسته‌ی برنامه‌ریزی $t \in \{0, 1, \ldots, T\}$ یک تصمیم
تکرارشونده می‌گیرد: بودجه‌ی پولی $x_{i,t} \geq 0$ که به هر پروژه‌ی
فعال $i$ داده می‌شود. افق اپیزود برابر است با:
\begin{equation}
  T = \max_{i} \bigl(f_i + \Delta^{\max}_i\bigr) - 1,
  \label{eq:horizon}
\end{equation}
که در آن $f_i = s_i + D_i$ پایان برنامه‌ریزی‌شده و $\Delta^{\max}_i$
حداکثر تأخیر مجاز در پایان پروژه‌ی $i$ است. دوره‌ی $t=0$ با زودترین
شروع برنامه‌ریزی‌شده‌ی پروژه‌ها هم‌تراز است. همه‌ی پیامدهای مورد توجه،
یعنی تحقق نقاط عطف، ورودی‌های نقدی، پایبندی به برنامه و فسخ پروژه، از
مسیر تصمیم تخصیص می‌گذرند. انتخاب سبد یک ورودی بیرونی در نظر گرفته
می‌شود~\citep{Archer1999}.

این صورت‌بندی بر چهار مفروضه‌ی ساختاری استوار است.

\begin{enumerate}[label=(\roman*)]
  \item \emph{قراردادهای دوجانبه‌ی مستقل.} پروژه‌ها فقط از راه موجودی
    نقدی مشترک سبد با هم ارتباط دارند. منابع یا نیروی کار مشترکی وجود
    ندارد و این با رویه‌ی استاندارد قراردادهای \lr{EPC} سازگار است.

  \item \emph{بودجه به‌عنوان تنها متغیر کنترلی.} تصمیم‌های داخلی
    پیمانکار درباره‌ی منابع در یک اسکالر کارایی تصادفی
    $\eta_{i,t} \geq 0.01$ خلاصه می‌شود. مدیر فقط می‌بیند که یک تخصیص
    بودجه چه مقدار پیشرفت فیزیکی می‌سازد، و نمی‌بیند این پیشرفت چگونه
    به‌دست آمده است. $\eta_{i,t}$ تنها متغیر تصادفی بنیادین مدل است و
    پیش از تعهد $x_{i,t}$ قابل مشاهده نیست.

  \item \emph{دوره‌های گسسته‌ی ماهانه.} واحد دوره یک ماه تقویمی است که
    چرخه‌ی استاندارد راهبری را بازتاب می‌دهد. نسبت تنگنای بودجه
    $\kappa = B_0 / \sum_i \mathrm{BAC}_i$ رژیم نقدینگی را پارامتری
    می‌کند.

  \item \emph{سازوکارهای قراردادی \lr{FIDIC}.} پیش‌پرداخت و بازیابی آن،
    درآمد نقاط عطف، کسر و بازگشت سپرده‌ی حسن انجام کار و تسویه‌ی فسخ
    همگی قابل پیکربندی‌اند و طیف کامل شرایط خصوصی رایج در عمل را پوشش
    می‌دهند.
\end{enumerate}

% cashflow diagram
\begin{figure}[htbp]
  \centering
  \includegraphics[width=\textwidth]{figures/cashflow_timeline.pdf}
  \caption[خط زمانی جریان نقدی]{خط زمانی جریان نقدی پروژه‌ی $i$ تحت
    سازوکارهای قراردادی \lr{FIDIC}. تبدیل بودجه به پیشرفت تابع کارایی
    تصادفی $\eta_{i,t}$ است.}
  \label{fig:cashflow_timeline}
\end{figure}

% ------------------------------------------------------------
\subsubsection{عناصر اولیه‌ی محیط}
\label{subsec:environment}
% ------------------------------------------------------------

% - - - - - - - - - - - - - - - - - - - - - - - - - - - - - -
\paragraph{دینامیک پیشرفت و مخارج}
\label{subsubsec:progress}
% - - - - - - - - - - - - - - - - - - - - - - - - - - - - - -

فرض کنید $\tau_i$ زمان سپری‌شده‌ی پروژه‌ی $i$ از شروع برنامه‌ریزی‌شده‌ی
آن، یعنی $s_i$، باشد. این مقدار در هر دوره‌ای که پروژه فعال است یک
واحد زیاد می‌شود. بودجه‌ی $x_{i,t}$ از راه اسکالر کارایی تصادفی
$\eta_{i,t}$ به پیشرفت فیزیکی تبدیل می‌شود:
\begin{equation}
  EV_{i,t}
    = \min\!\Bigl(1,\;
        EV_{i,t-1} + \eta_{i,t} \cdot \frac{x_{i,t}}{\mathrm{BAC}_i}
      \Bigr),
  \quad
  \mathrm{ACWP}_{i,t} = \mathrm{ACWP}_{i,t-1} + x_{i,t},
  \label{eq:progress_rec}
\end{equation}
با شرط اولیه‌ی $EV_{i,0} = \mathrm{ACWP}_{i,0} = 0$. در این رابطه
$\mathrm{BAC}_i$ بودجه در زمان تکمیل است و $EV_{i,t} \in [0,1]$ کسر
کار انجام‌شده است (هزینه‌ی بودجه‌ای کار انجام‌شده تقسیم بر
$\mathrm{BAC}_i$). افزایش پیشرفت $\eta_{i,t} x_{i,t} /
\mathrm{BAC}_i$ فقط بعد از تعهد تخصیص محقق می‌شود. مقدار $\eta_{i,t}$
از توزیعی مخصوص هر پروژه نمونه‌برداری می‌شود که خانواده و پارامترهای
آن در دیکشنری پیکربندی مشخص شده‌اند.

پیشرفت تجمعی \emph{برنامه‌ریزی‌شده} از یک منحنی \lr{S} نرمال‌شده با
پارامترهای شکل بتا $(a_i, b_i)$ پیروی می‌کند~\citep{Kenley2009}:
\begin{equation}
  PV_i(\tau) = I_{\tau / D_i}(a_i,\, b_i),
  \quad \tau \in [0, D_i],\quad D_i \geq 1,
  \label{eq:scurve}
\end{equation}
که در آن $I_u(\cdot,\cdot)$ تابع بتای ناقص نرمال‌شده و
$D_i \in \mathbb{Z}_{\geq 1}$ مدت برنامه‌ریزی‌شده بر حسب دوره است.
کمترین مدت مجاز یک دوره است ($D_i = 1$). این حالت پروژه‌ای را مدل
می‌کند که کل دامنه‌ی برنامه‌ریزی‌شده‌ی آن باید با یک تخصیص واحد در
$t = s_i$ تحویل داده شود. تابع توزیع تجمعی بتا مقدار $\tau = 0$ را به
$0$ و $\tau = 1$ را به $1$ می‌برد، پس مدیر دقیقاً یک دوره فرصت دارد تا
تخصیص کامل را پیش از فعال‌شدن سررسید قطعی در $t = s_i + 1$ متعهد شود.
مقدار $PV_i(\tau_i)$ پیشرفت برنامه‌ریزی‌شده در زمان جاری پروژه است.
مقادیر $PV_i(\tau_i + 1)$ و $PV_i(\tau_i + 2)$ به‌ترتیب نگاه یک‌دوره‌ای
و دودوره‌ای به آینده‌اند و در بردار مشاهده و سیاست‌های خط مبنا به‌کار
می‌روند.

% - - - - - - - - - - - - - - - - - - - - - - - - - - - - - -
\paragraph{سازوکارهای جریان نقدی \lr{FIDIC}}
\label{subsubsec:cashflow}
% - - - - - - - - - - - - - - - - - - - - - - - - - - - - - -

\subparagraph{قیمت قرارداد و پیش‌پرداخت.}
قیمت قرارداد برابر $P_i = \mathrm{BAC}_i (1 + \pi_i)$ است، که در آن
$\pi_i$ حاشیه‌ی سود است. اگر قرارداد پیش‌پرداخت داشته باشد
($\phi_i > 0$)، سبد در شروع پروژه ($t = s_i$) مبلغ زیر را دریافت
می‌کند:
\begin{equation}
  \mathrm{Adv}_i = \phi_i \cdot P_i,
  \label{eq:advance}
\end{equation}
که در آن $\phi_i$ نرخ پیش‌پرداخت است. وقتی $\phi_i = 0$ باشد، هیچ
پیش‌پرداختی پرداخت نمی‌شود، سازوکار بازیابی غیرفعال می‌ماند و برای همه‌ی
نقاط عطف $A^{\mathrm{rec}}_{i,j} = 0$ است. پیش‌پرداخت به‌شکل یک نقطه‌ی
عطف مصنوعی با شماره‌ی $j = 0$ و آستانه‌ی پیشرفت صفر مدل می‌شود که
به‌طور خودکار در $t = s_i$ تأیید می‌شود. به این ترتیب
$\hat{P}_{i,0} = \mathrm{Adv}_i$ بدون نیاز به حالت خاص، از راه
رابطه‌ی~\eqref{eq:Vt} وارد $V_{s_i}$ می‌شود.

بازیابی پیش‌پرداخت از پرداخت‌های بعدی نقاط عطف، از سازوکار بند فرعی
۱۴٫۲ \lr{FIDIC} پیروی می‌کند~\citep{FIDIC2017}. این سازوکار با دو
پارامتر مخصوص هر پروژه اداره می‌شود: آستانه‌ی شروع
$\psi^{\mathrm{rec}}_i \in [0,1]$ و نرخ کسر از هر پرداخت
$\delta^{\mathrm{rec}}_i \in [0,1]$. فرض کنید
$G^{\mathrm{cum}}_{i,j} = \sum_{k=1}^{j} G_{i,k}$ مجموع تجمعی پرداخت
ناخالص تأییدشده تا نقطه‌ی عطف $j$ باشد. مبلغ بازیابی پیش‌پرداخت در
نقطه‌ی عطف $j$ چنین است:
\begin{equation}
  A^{\mathrm{rec}}_{i,j}
    = \begin{cases}
        \min\!\bigl(\delta^{\mathrm{rec}}_i \cdot G_{i,j},\;
          \mathrm{Adv}^{\mathrm{rem}}_{i,j}\bigr)
          & \text{اگر } G^{\mathrm{cum}}_{i,j-1} \geq
            \psi^{\mathrm{rec}}_i \cdot P_i
            \text{ و } \phi_i > 0, \\
        0 & \text{در غیر این صورت،}
      \end{cases}
  \label{eq:advance_rec}
\end{equation}
که در آن $\mathrm{Adv}^{\mathrm{rem}}_{i,j}$ مانده‌ی پیش‌پرداخت
تسویه‌نشده پیش از تأیید نقطه‌ی عطف $j$ است و تضمین می‌کند مجموع مبلغ
بازیابی‌شده هرگز از $\mathrm{Adv}_i$ بیشتر نشود. وقتی
$\psi^{\mathrm{rec}}_i = 0$ باشد، کسر بازیابی از همان اولین پرداخت
نقطه‌ی عطف شروع می‌شود.

\subparagraph{تأیید نقطه‌ی عطف و پرداخت‌های خالص.}
نقاط عطف $j = 1, \ldots, M_i$ بر اساس آستانه‌ی پیشرفت
$\psi_{i,1} < \cdots < \psi_{i,M_i} = 1$ مرتب شده‌اند. نقطه‌ی عطف $j$
در نخستین دوره‌ی $t$ تأیید می‌شود که هم‌زمان $EV_{i,t} \geq \psi_{i,j}$
و $t \geq s_i + \tau^{\mathrm{ms}}_{i,j}$ برقرار باشد. در اینجا
$\tau^{\mathrm{ms}}_{i,j}$ آستانه‌ی زمانی پروژه است که فاصله‌ی پرداخت
مندرج در بند فرعی ۱۴٫۳ \lr{FIDIC} را اعمال می‌کند~\citep{FIDIC2017}.
تأیید یکنواخت است. وجه نقد خالصی که هنگام تأیید نقطه‌ی عطف $j$ دریافت
می‌شود برابر است با:
\begin{equation}
  \hat{P}_{i,j}
    = G_{i,j}
      - A^{\mathrm{rec}}_{i,j}
      - R^{\mathrm{held}}_{i,j}
      + R^{\mathrm{rel}}_{i,j},
  \label{eq:net_payment}
\end{equation}
که در آن $G_{i,j} = w_{i,j} P_i$ پرداخت ناخالص است ($w_{i,j}$ وزن
نقطه‌ی عطف با $\sum_{j=1}^{M_i} w_{i,j} = 1$)،
$A^{\mathrm{rec}}_{i,j}$ مبلغ بازیابی پیش‌پرداخت
در رابطه‌ی~\eqref{eq:advance_rec} است،
$R^{\mathrm{held}}_{i,j} = \nu_i G_{i,j}$ سپرده‌ی حسن انجام کار
نگه‌داشته‌شده است، و $R^{\mathrm{rel}}_{i,j}$ فقط در آخرین نقطه‌ی عطف
برابر مجموع تجمعی سپرده‌ی حسن انجام کار است و در بقیه‌ی نقاط صفر است.
وقتی $\phi_i > 0$ باشد، اتحاد کامل پرداخت
$\mathrm{Adv}_i + \sum_{j=1}^{M_i} \hat{P}_{i,j} = P_i$ دقیقاً برقرار
است. وقتی $\phi_i = 0$ باشد، این اتحاد به
$\sum_{j=1}^{M_i} \hat{P}_{i,j} = P_i$ تبدیل می‌شود. پرداخت‌ها در همان
دوره‌ی تأیید دریافت می‌شوند.

\subparagraph{ورودی تجمعی دوره.}
کل ورودی نقدی سبد در دوره‌ی $t$ برابر است با:
\begin{equation}
  V_t
    = \sum_{i=0}^{N-1} \sum_{j=0}^{M_i}
        \hat{P}_{i,j} \cdot
        \mathbf{1}\!\bigl[t^{\mathrm{cert}}_{i,j} = t\bigr]
      + \sum_{i \in \mathcal{P}_t^{\mathrm{term}}} S_i,
  \label{eq:Vt}
\end{equation}
که در آن $t^{\mathrm{cert}}_{i,j}$ دوره‌ی تأیید نقطه‌ی عطف $j$ است
(اگر تأیید نشده باشد، تهی)، $\mathcal{P}_t^{\mathrm{term}}$ مجموعه‌ی
پروژه‌هایی است که در $t$ تازه فسخ شده‌اند، و $S_i$ تسویه‌ی فسخ است که
در بخش~\ref{subsubsec:termination} تعریف می‌شود.

% - - - - - - - - - - - - - - - - - - - - - - - - - - - - - -
\paragraph{سیگنال‌های عملکرد \lr{EVM}}
\label{subsubsec:evm}
% - - - - - - - - - - - - - - - - - - - - - - - - - - - - - -

چارچوب مدیریت ارزش کسب‌شده~\citep{Fleming2010} چهار شاخص تشخیصی در
اختیار می‌گذارد که در هر دوره از حالت قابل مشاهده محاسبه می‌شوند.
زمان کسب‌شده‌ی $\mathrm{ES}_{i,t}$ با وارون‌کردن منحنی \lr{S}
برنامه‌ریزی‌شده در سطح پیشرفت واقعی فعلی به‌دست می‌آید. یعنی
$\mathrm{ES}_{i,t}$ همان زمان پروژه‌ی $\tau^* \in [0, D_i]$ است که
$I_{\tau^*/D_i}(a_i, b_i) = EV_{i,t}$ شود و با جست‌وجوی دودویی محاسبه
می‌شود. چهار شاخص عبارت‌اند از:
\begin{align}
  \mathrm{SPI}_{i,t}
    &= \frac{\mathrm{ES}_{i,t}}{\tau_i},
  \label{eq:SPI} \\
  \mathrm{CPI}_{i,t}
    &= \frac{EV_{i,t} \cdot \mathrm{BAC}_i}{\mathrm{ACWP}_{i,t}},
  \label{eq:CPI} \\
  \mathrm{EAC}_{i,t}
    &= \mathrm{ACWP}_{i,t}
       + \frac{\mathrm{BAC}_i - EV_{i,t} \cdot \mathrm{BAC}_i}
              {\mathrm{CPI}_{i,t} \cdot \mathrm{SPI}_{i,t}},
  \label{eq:EAC} \\
  \mathrm{TCPI}_{i,t}
    &= \frac{\mathrm{BAC}_i - EV_{i,t} \cdot \mathrm{BAC}_i}
            {\mathrm{BAC}_i - \mathrm{ACWP}_{i,t}}.
  \label{eq:TCPI}
\end{align}

شاخص عملکرد زمان‌بندی مبتنی بر زمان کسب‌شده~\eqref{eq:SPI} به‌جای نسبت
کلاسیک $\mathrm{SPI} = \mathrm{EV}/\mathrm{PV}$ به‌کار می‌رود، چون پس
از تاریخ پایان برنامه‌ریزی‌شده هم به کاهش خود ادامه می‌دهد و پروژه‌ای
را که یک دوره دیرکرد دارد از پروژه‌ای که چند دوره دیرکرد دارد به‌درستی
جدا می‌کند~\citep{Lipke2003}. مخرج ترکیبی در رابطه‌ی~\eqref{eq:EAC}
ناکارآمدی هزینه و زمان‌بندی را هر دو در پیش‌بینی هزینه‌ی در زمان تکمیل
وارد می‌کند~\citep{christensen1993est}. مقدار $\mathrm{TCPI} > 1.10$
به‌طور گسترده مسیر غیرقابل جبران افزایش هزینه به‌شمار می‌آید~\citep{Fleming2010}.
مقادیر $\mathrm{SPI}_{i,t}$، $\mathrm{CPI}_{i,t}$ و
$\mathrm{EAC}_{i,t}$ در دفتر ثبت حسابرسی ذخیره می‌شوند. این مقادیر
متغیرهای میانی‌اند که ویژگی‌های تخصیص جبرانی و نقطه‌ی عطف هدف در بردار
مشاهده از آن‌ها ساخته می‌شود، اما خودشان در بردار مشاهده‌ی $o_t$ قرار
نمی‌گیرند. مخرج‌هایی که عملاً صفرند به‌طور پیش‌فرض برابر یک گرفته
می‌شوند.

% - - - - - - - - - - - - - - - - - - - - - - - - - - - - - -
\paragraph{سازوکار فسخ}
\label{subsubsec:termination}
% - - - - - - - - - - - - - - - - - - - - - - - - - - - - - -

فسخ یک سازوکار درون‌زای محیط است و عمل عامل نیست. در هر دوره دو شرط
به‌طور مستقل بررسی می‌شود.

\subparagraph{شرط ۱ --- نقض ناشی از تأخیر پیشرفت.}
تأخیر علامت‌دار پیشرفت در دوره‌ی $t$ برابر است با:
\begin{equation}
  \delta_{i,t} = PV_i(\tau_i) - EV_{i,t}.
  \label{eq:progress_lag}
\end{equation}
وقتی تأخیر از سقف قراردادی بیشتر شود، پرچم نقض فعال می‌شود:
\begin{equation}
  \mathbb{1}^{\delta}_{i,t}
    = \mathbf{1}\!\bigl[\delta_{i,t} > \delta^{\max}_i\bigr],
  \label{eq:term_cond1}
\end{equation}
که در آن $\delta^{\max}_i \in [0,1]$ در زمان تعهد ثابت می‌شود و مبنای
آن اختیار مهندس \lr{FIDIC} برای صدور اخطار در زمانی است که پیشرفت
به‌طور محسوسی از برنامه عقب بیفتد~\citep{FIDIC2017}.

\subparagraph{شرط ۲ --- سررسید قطعی.}
سررسید قطعی بدون هیچ قید و شرطی در دوره‌ی زیر فعال می‌شود:
\begin{equation}
  \mathbb{1}^{T}_{i,t}
  = \mathbf{1}\!\bigl[t = f_i + \Delta^{\max}_i - 1\bigr],
  \label{eq:term_cond2}
\end{equation}
و فسخ فوری را، صرف‌نظر از وضعیت شمارنده‌ی مهلت جبران، به‌دنبال دارد.

\subparagraph{شمارنده‌ی مهلت جبران.}
شرط ۱ یک شمارنده‌ی مهلت جبران را فعال می‌کند که فاصله‌ی جبران در بند
فرعی ۱۵٫۱ \lr{FIDIC} را مدل می‌کند~\citep{FIDIC2017, Ndekugri2008}:
\begin{equation}
  \theta^{\mathrm{rem}}_{i,t}
    = \begin{cases}
        \max\!\bigl(0,\; \theta^{\mathrm{rem}}_{i,t-1} - 1\bigr)
          & \text{اگر } \mathbb{1}^{\delta}_{i,t} = 1, \\
        \theta_i
          & \text{در غیر این صورت،}
      \end{cases}
  \label{eq:counter}
\end{equation}
با مقدار اولیه‌ی $\theta^{\mathrm{rem}}_{i,s_i} = \theta_i$. فسخ نرم
وقتی انجام می‌شود که $\theta^{\mathrm{rem}}_{i,t} = 0$ شود. فسخ
قطعی~\eqref{eq:term_cond2} در فاز آغازین (پیش از تخصیص) بررسی می‌شود،
و شرط ۱ و به‌روزرسانی شمارنده در فاز پایانی (پس از تخصیص و تحقق
پیشرفت) اعمال می‌شوند. شمارنده‌ی $\theta^{\mathrm{rem}}_{i,t}$ جزو
حالت قابل مشاهده است و مستقیماً در بردار مشاهده می‌آید. عامل به‌طور
غیرمستقیم و با تخصیص‌هایی که $EV_{i,t}$ را نسبت به $PV_i(\tau_i)$
جابه‌جا می‌کنند روی آن اثر می‌گذارد.

\subparagraph{تسویه‌ی فسخ.}
وقتی پروژه‌ی $i$ در دوره‌ی $t$ فسخ شود، جریان نقدی خالص تسویه برابر
است با:
\begin{equation}
  S_i = EV_{i,t} \cdot P_i - \mathrm{CF}^+_i,
  \label{eq:settlement}
\end{equation}
که در آن $\mathrm{CF}^+_i$ ورودی تجمعی دریافت‌شده برای پروژه‌ی $i$ تا
دوره‌ی $t$ است. جمله‌ی اول حق پیمانکار بابت کار فیزیکی تحویل‌داده‌شده
است و جمله‌ی دوم آن چیزی است که پیش‌تر دریافت شده است. وقتی کار
تحویل‌شده از دریافتی‌های قبلی بیشتر باشد، $S_i > 0$ است. اگر
$S_i < 0$ باشد، یک تعهد خالص بر دوش سبد می‌افتد (پیش‌پرداخت بازیابی‌نشده).
در هر دو حالت، $S_i$ کاملاً با پارامترهای قرارداد و سطح پیشرفت در زمان
فسخ تعیین می‌شود و به کالیبراسیون بیرونی نیازی ندارد.

ماشین حالت حاکم بر چرخه‌ی عمر هر پروژه در
شکل~\ref{fig:state_machine} نشان داده شده است.

\begin{figure}[H]
  \centering
  \includegraphics[width=\textwidth]{figures/state_machine.pdf}
  \caption[منطق حالت فسخ]{پویایی شمارنده‌ی تحمل نقض: نقض ناشی از تأخیر
    شمارنده را کم می‌کند و عبور از سررسید قطعی فسخ فوری را به‌دنبال
    دارد. تسویه‌ی $S_i$ \eqref{eq:settlement} هنگام فسخ وارد ورودی‌های
    سبد می‌شود.}
  \label{fig:state_machine}
\end{figure}

% - - - - - - - - - - - - - - - - - - - - - - - - - - - - - -
\paragraph{موجودی خزانه و بهره}
\label{subsubsec:balance}
% - - - - - - - - - - - - - - - - - - - - - - - - - - - - - -

جریان نقدی خالص در دوره‌ی $t$ برابر
$\Delta B_t = V_t - \sum_{i \in \mathcal{P}_t^{\mathrm{active}}} x_{i,t}
- \sum_i c^r_{i,t}$ است، که در آن $c^r_{i,t}$ بهره‌ی برداشت از خزانه
است که به پروژه‌ی $i$ تعلق می‌گیرد:
\begin{equation}
  c^r_{i,t}
    = \frac{r}{12} \cdot
      \max\!\bigl(0,\; \mathrm{CF}^-_{i,t-1} + x_{i,t}
        - \mathrm{CF}^+_{i,t-1}\bigr),
  \label{eq:interest}
\end{equation}
که در آن $\mathrm{CF}^-_{i,t}$ و $\mathrm{CF}^+_{i,t}$ به‌ترتیب
خروجی‌ها و ورودی‌های تجمعی پروژه‌ی $i$ تا دوره‌ی $t$ هستند و $r$ نرخ
بهره‌ی سالانه است. بهره وقتی انباشته می‌شود که مخارج تجمعی یک پروژه از
دریافتی‌های تجمعی آن بیشتر شود، یعنی وقتی پروژه به‌جای تأمین مالی از
ورودی‌های نقاط عطف، از ذخایر خزانه برداشت می‌کند. موجودی سبد به این
شکل تحول می‌یابد:
\begin{equation}
  B_{t+1} = B_t + \Delta B_t.
  \label{eq:balance}
\end{equation}
وقتی $B_t \leq 0$ باشد و هیچ پیش‌پرداخت در انتظاری وجود نداشته باشد،
اپیزود بلافاصله بریده می‌شود.

% ------------------------------------------------------------
\subsubsection{مرجع پارامترها}
\label{subsec:parameters}
% ------------------------------------------------------------

جدول~\ref{tab:params} همه‌ی پارامترهایی را که در زمان تعهد ثابت
می‌شوند یک‌جا گردآوری می‌کند.

\begin{table}[H]
  \centering
  \includegraphics[width=\textwidth]{tables/params.pdf}
  \caption[جدول پارامترها]{پارامترهای مدل و دامنه‌ی آن‌ها.}
  \label{tab:params}
\end{table}

% ------------------------------------------------------------
\subsubsection{کران بالای اطلاعات‌کامل و شدنی‌نبودن محاسباتی}
\label{subsec:theoretical_ub}
% ------------------------------------------------------------

\paragraph{کران بالای اوراکل قطعی}
\label{subsubsec:ub}

یک اوراکل فرضی را در نظر بگیرید که پیش از گرفتن هر تصمیمی، کل دنباله‌ی
کارایی $(\eta_{i,t})_{i,t}$ را می‌داند. وقتی کارایی قطعی باشد،
$EV_{i,t}$ تابعی قطعی از $\{x_{i,\tau}\}_{\tau \leq t}$ است، پس زمان
تأیید همه‌ی نقاط عطف و دوره‌های فسخ از پیش قابل محاسبه‌اند. اوراکل
مسئله‌ی زیر را بیشینه می‌کند:
\begin{equation}
  Z^* = \max_{\{x_{i,t} \geq 0\}}
    \sum_{t=0}^{T} \gamma^{t} \Delta B_t,
  \label{eq:ub_obj}
\end{equation}
با قید $\sum_{i \in \mathcal{P}_t^{\mathrm{active}}} x_{i,t}
\leq B_t$ برای همه‌ی $t$، در حالی که موجودی طبق
رابطه‌ی~\eqref{eq:balance} تحول می‌یابد و همه‌ی ورودی‌های $V_t$ به‌طور
قطعی از مسیر کارایی معلوم محاسبه می‌شوند.

\begin{proposition}[کران بالای اطلاعات‌کامل]
\label{prop:upper_bound}
فرض کنید $Z^*$ مقدار رابطه‌ی~\eqref{eq:ub_obj} روی یک مسیر تحقق‌یافته‌ی
کارایی باشد. فرض کنید $Z^\pi$ پاداش تجمعی تنزیل‌شده‌ای باشد که هر
سیاست علّی $\pi$ روی همان مسیر به‌دست می‌آورد؛ سیاست علّی سیاستی است
که $x_{i,t}$ را فقط با اطلاعات در دسترس در زمان $t$ تعهد می‌کند. در
این صورت:
\begin{equation}
  \sup_{\pi \in \Pi^{\mathrm{causal}}}
    \mathbb{E}\!\bigl[Z^\pi\bigr]
  \leq Z^*.
  \label{eq:upper_bound}
\end{equation}
\end{proposition}

\begin{proof}
یک مسیر تحقق‌یافته‌ی دلخواه را ثابت بگیرید. هر دنباله‌ی تخصیصی که یک
سیاست علّی تولید کند، در مسئله‌ی~\eqref{eq:ub_obj} که روی همان
تحقق‌ها حل می‌شود هم شدنی است، چون قیدهای شدنی‌بودن یکسان‌اند و
اوراکل روی مجموعه‌ی اطلاعاتی بسیار بزرگ‌تری بهینه‌سازی می‌کند. پس
به‌ازای هر مسیر $Z^\pi \leq Z^*$ برقرار است و این نامساوی در امید
ریاضی هم حفظ می‌شود.
\end{proof}

شکاف $Z^* - \mathbb{E}[Z^\pi]$ به هزینه‌ی عدم‌قطعیت و کیفیت سیاست
تجزیه می‌شود. در این پژوهش، $Z^*$ برای هر پیکربندی آزمون به‌صورت دستی
و با آگاهی کامل از نمونه‌های کارایی تحقق‌یافته در هر جلسه‌ی بازی،
از راه رابط تعاملی \texttt{play.py} محاسبه می‌شود. این جواب‌های
دستی مرجع کران بالا را می‌سازند که در فصل~\ref{chap:evaluation}
گزارش شده است.

\paragraph{شدنی‌نبودن بهینه‌سازی تصادفی دقیق}
\label{subsubsec:intractability}

همتای تصادفی دقیق رابطه‌ی~\eqref{eq:ub_obj} باید روی سیاست‌هایی بهینه
شود که تاریخچه‌ی مشاهده‌شده را به تخصیص‌ها نگاشت می‌کنند و هم‌زمان
قید عدم‌پیش‌نگری را رعایت می‌کنند، یعنی $x_{i,t}$ فقط به اطلاعات در
دسترس در زمان $t$ وابسته باشد. روش حل متداول برای این کار شمارش درخت
سناریوهاست~\citep{Birge2011}. اگر $\eta_{i,t}$ را برای هر پروژه و هر
دوره به $K$ مقدار گسسته کنیم، درخت تعداد برگ زیر را دارد:
\begin{equation}
  |\mathcal{S}| = K^{N \cdot T}.
  \label{eq:tree_size}
\end{equation}
برای یک سبد متوسط با $N = 3$ پروژه، $T = 24$ دوره و $K = 5$ سناریوی
گسسته برای هر پروژه در هر دوره، داریم
$|\mathcal{S}| = 5^{72} \approx 10^{50}$. با $N = 10$ پروژه، این عدد
به $5^{240} \approx 10^{167}$ می‌رسد که بیش از ۸۰ مرتبه‌ی بزرگی از
تعداد اتم‌های جهان قابل مشاهده بیشتر است. این رشد از ذات صورت‌بندی
تصادفی چندمرحله‌ای می‌آید و به فناوری حل‌کننده بستگی ندارد. دینامیک
درون‌زای مجموعه‌ی فعال منبع دومی از رشد ترکیباتی است که بیش از
رابطه‌ی~\eqref{eq:tree_size} است: هر رویداد فسخ، مجموعه‌ی
$\mathcal{P}_t^{\mathrm{active}}$ را ساختاری عوض می‌کند و درخت
سناریو را شاخه‌شاخه می‌کند.

این شدنی‌نبودن، صورت‌بندی \lr{MDP} در بخش~\ref{subsec:mdp} را توجیه
می‌کند: شمارش درخت سناریو با تجربه‌ی مسیرهای نمونه‌برداری‌شده جایگزین
می‌شود و هزینه‌ی محاسباتی هر گام به $|\mathcal{S}|$ بستگی ندارد.

% ------------------------------------------------------------
\subsubsection{فرمول‌بندی فرایند تصمیم مارکوف}
\label{subsec:mdp}
% ------------------------------------------------------------

% - - - - - - - - - - - - - - - - - - - - - - - - - - - - - -
\paragraph{تعریف رسمی \lr{MDP}}
\label{subsubsec:mdp_def}
% - - - - - - - - - - - - - - - - - - - - - - - - - - - - - -

\begin{definition}[فرایند تصمیم مارکوف تخصیص بودجه‌ی سبد]
\label{def:mdp}
مسئله‌ی تخصیص بودجه در سبد به‌صورت یک فرایند تصمیم مارکوف با افق
متناهی رسمی می‌شود~\citep{Puterman1994}:
\begin{equation}
  \mathcal{M} = \bigl(\mathcal{S},\;\mathcal{A},\;
    \mathcal{P},\;R,\;\gamma,\;T\bigr),
  \label{eq:mdp_tuple}
\end{equation}
با فضای حالت $\mathcal{S}$، فضای عمل $\mathcal{A}$، هسته‌ی گذار
$\mathcal{P}$، تابع پاداش $R$، ضریب تنزیل $\gamma$ و افق $T$.
\end{definition}

\subparagraph{فضای حالت.}
حالت $S_t$ در دوره‌ی $t$ برای هر پروژه‌ی $i$ مؤلفه‌های زیر را در بر
می‌گیرد:
\begin{itemize}[nosep]
  \item پیشرفت واقعی $EV_{i,t} \in [0,1]$،
  \item هزینه‌ی واقعی تجمعی $\mathrm{ACWP}_{i,t} \in \mathbb{R}_{\geq 0}$،
  \item شاخص‌های \lr{EVM} یعنی $\mathrm{SPI}_{i,t},\; \mathrm{CPI}_{i,t},\;
        \mathrm{EAC}_{i,t},\; \mathrm{TCPI}_{i,t}$،
  \item وضعیت پروژه‌ی $\sigma_{i,t} \in \{\texttt{pending},\,
        \texttt{active},\, \texttt{completed},\, \texttt{terminated}\}$،
  \item شمارنده‌ی مهلت جبران $\theta^{\mathrm{rem}}_{i,t} \in
        \{0, \ldots, \theta_i\}$،
\end{itemize}
به‌همراه موجودی نقدی سبد $B_t \in \mathbb{R}_{\geq 0}$. حالت کامل
چنین است:
\begin{equation}
  S_t \;\in\; \mathcal{S}
  \;\subseteq\; [0,1]^N \times \mathbb{R}^N_{\geq 0}
    \times \mathbb{R}^{4N} \times \Sigma^N
    \times \mathbb{Z}^N_{\geq 0} \times \mathbb{R}_{\geq 0},
  \label{eq:state_space}
\end{equation}
که در آن $\Sigma = \{\texttt{pending},\texttt{active},\texttt{completed},
\texttt{terminated}\}$. بُعد حالت در هر دوره برابر $9N+1$ است. همه‌ی
مؤلفه‌ها در زمان تصمیم‌گیری قابل مشاهده‌اند. تنها کمیت مشاهده‌ناپذیر
$\eta_{i,t}$ است که فقط بعد از تعهد $x_{i,t}$ محقق می‌شود.

خاصیت مارکوف دقیقاً برقرار است: $S_{t+1}$ فقط به
$(S_t, \mathbf{x}_t, \boldsymbol{\eta}_t)$ بستگی دارد و به تاریخچه‌ی
قبلی $(S_0, \mathbf{x}_0, \ldots, S_{t-1}, \mathbf{x}_{t-1})$ وابسته
نیست. حالت کامل سبد $S_t$ یک آماره‌ی بسنده برای تاریخچه‌ی قابل مشاهده
است، پس صورت‌بندی \lr{MDP} دقیق است و تقریبی نیست.

\subparagraph{فضای عمل.}
سیاست یک بردار کسر تخصیص $\boldsymbol{\alpha}_t \in [0,1]^N$ تولید
می‌کند، با این تفسیر که پروژه‌ی $i$ کسر $\alpha_{i,t}$ از بودجه‌ی در
دسترس را دریافت می‌کند. تخصیص پولی برابر
$x_{i,t} = \alpha_{i,t} \cdot B_t$ است.

\subparagraph{تصویرسازی عمل.}
برای اینکه هر خروجی خام سیاست،
$\tilde{\boldsymbol{\alpha}}_t \in \mathbb{R}^N$، شدنی شود، محیط پیش
از اجرای گذار سه قاعده‌ی تصویرسازی را پشت سر هم اعمال می‌کند:

\begin{enumerate}[label=(\roman*)]
  \item \emph{صفرکردن پروژه‌های غیرفعال.}
    $\tilde{\alpha}_{i,t} \leftarrow 0$ برای همه‌ی $i$ با
    $\sigma_{i,t} \neq \texttt{active}$.

  \item \emph{سقف تخصیص برای هر پروژه.}
    تخصیص پولی پروژه‌ی $i$ به مقدار زیر محدود می‌شود:
    \begin{equation}
      x_{i,t} \leq \frac{\mathrm{BAC}_i}{\eta^{\min}_i},
      \label{eq:alloc_cap}
    \end{equation}
    که در آن $\eta^{\min}_i$ کران پایین توزیع کارایی پروژه‌ی $i$ است
    (برای خانواده‌های ثابت و یکنواخت به‌صورت فرمول بسته، و در غیر این
    صورت مقدار پیش‌فرض ۰٫۰۱). این سقف تضمین می‌کند که یک تخصیص کامل
    واحد، هر نمونه‌ی کارایی که باشد، همیشه برای تکمیل پروژه کافی است.

  \item \emph{تصویرسازی بودجه.}
    زیربردار فعال روی سیمپلکس بودجه تصویر می‌شود:
    \begin{equation}
      \boldsymbol{\alpha}_t^{\mathrm{active}}
        = \frac{\max(\tilde{\boldsymbol{\alpha}}_t^{\mathrm{active}},\,\mathbf{0})}
               {\bigl\|\max(\tilde{\boldsymbol{\alpha}}_t^{\mathrm{active}},\,
                 \mathbf{0})\bigr\|_1 \vee 1},
      \label{eq:projection}
    \end{equation}
    که در آن $\vee$ عملگر بیشینه است (وقتی همه‌ی مقادیر نامثبت باشند،
    مخرج دست‌کم یک می‌شود). تخصیص‌های پولی حاصل، یعنی
    $x_{i,t} = \alpha_{i,t} B_t$، طبق ساختار
    $\sum_i x_{i,t} \leq B_t$ را برآورده می‌کنند.
\end{enumerate}

این تصویرسازی تقریباً همه‌جا مشتق‌پذیر است و در زمان آموزش، جریان
گرادیان را در شبکه‌های سیاست حفظ می‌کند~\citep{Schulman2017}.

\subparagraph{هسته‌ی گذار.}
گذار تصادفی $\mathcal{P}(S_{t+1} \mid S_t, \boldsymbol{\alpha}_t)$ با
الگوریتم~\ref{alg:kernel} اجرا می‌شود.

\begin{algorithm}[H]
  \centering
  \includegraphics[width=\textwidth]{algorithms/alg_kernel.pdf}
  \caption[هسته‌ی گذار]{هسته‌ی گذار درون‌دوره‌ای در گام $t$ از اپیزود.
    ردیف‌ها به‌ترتیب شماره‌گذاری شده‌اند و تورفتگی عمق تودرتویی را
    نشان می‌دهد. ستون راست شماره‌ی رابطه‌ی مربوط به هر قاعده‌ی
    به‌روزرسانی را می‌دهد.}
  \label{alg:kernel}
\end{algorithm}

% - - - - - - - - - - - - - - - - - - - - - - - - - - - - - -
\paragraph{تابع پاداش}
\label{subsubsec:reward}
% - - - - - - - - - - - - - - - - - - - - - - - - - - - - - -

پاداش هر دوره، جریان نقدی خالص تنزیل‌شده‌ی سبد است:
\begin{equation}
  r_t = \gamma^{t} \cdot \Delta B_t
      = \gamma^{t} \!\left(
          V_t
          - \sum_{i \in \mathcal{P}_t^{\mathrm{active}}} x_{i,t}
          - \sum_{i} c^r_{i,t}
        \right).
  \label{eq:reward}
\end{equation}
پیش‌پرداخت‌هایی که در $t = s_i$ برای پروژه‌های با $s_i = 0$ دریافت
می‌شوند، هنگام بازنشانی در بودجه‌ی اولیه‌ی $B_0$ جذب می‌شوند و از
$\Delta B_t$ بیرون می‌مانند. پیش‌پرداخت‌ها در شروع قرارداد و بدون
توجه به هیچ تصمیم تخصیصی، بی‌قیدوشرط پرداخت می‌شوند. اگر آن‌ها را در
$r_0$ بیاوریم، بازده هر اپیزود به اندازه‌ی یک مقدار ثابت بالا می‌رود
که هیچ اطلاعاتی درباره‌ی کیفیت سیاست ندارد. این نوعی سوگیری پاداش است
که سیگنال واقعی عملکرد را می‌پوشاند، بی‌آنکه رتبه‌بندی سیاست بهینه را
عوض کند.

هدف بهینه‌سازی سیاست به‌صورت زیر است:
\begin{equation}
  \pi^* = \argmax_{\pi}
    \;\mathbb{E}_{\pi}\!\left[
      G = \sum_{t=0}^{T} r_t
    \right],
  \label{eq:mdp_objective}
\end{equation}
که در آن امید ریاضی روی توزیع مسیرهای ناشی از $\pi$ و نمونه‌های
کارایی $\{\eta_{i,t}\}$ گرفته می‌شود.

\subparagraph{ویژگی‌های انگیزشی ساختاری.}
پاداش~\eqref{eq:reward} همه‌ی انگیزه‌های عملکردی را از راه ساختار مالی
محیط و بدون جمله‌های جریمه‌ی کمکی در خود دارد:

\begin{enumerate}[label=(\roman*)]
  \item \emph{انگیزه‌ی تکمیل.} هنگام تکمیل پروژه، سپرده‌ی حسن انجام
    کار انباشته‌شده‌ی $\nu_i P_i$ همراه با پرداخت آخرین نقطه‌ی عطف
    آزاد می‌شود. این یک ورودی نقدی متمرکز می‌سازد که سیاست‌هایی را
    پاداش می‌دهد که پروژه‌ها را تا تأیید کامل پیش می‌برند.

  \item \emph{بازدارندگی فسخ.} تسویه‌ی منفی $S_i$ به‌صورت کاهش مستقیم
    درآمد وارد $V_t$ می‌شود و بدون هیچ ضریب جریمه‌ی دست‌ساز، فسخ را
    از نظر مالی پرهزینه می‌کند. شدت این تهدید با فاصله‌ی میان
    پیش‌پرداخت‌های پرداخت‌شده و ارزش تأییدشده‌ی نقاط عطف متناسب است
    و عامل را تشویق می‌کند تخصیص‌ها را با پیشرفت واقعی هماهنگ کند.

  \item \emph{انگیزه‌ی زمان‌بندی.} وقتی $\gamma < 1$ باشد، تأیید زودتر
    نقطه‌ی عطف $j$ باعث می‌شود $\hat{P}_{i,j}$ با وزن تنزیل بیشتری
    دریافت شود. بنابراین سیاستی که بازده را بیشینه می‌کند، تخصیص‌ها را
    به‌سمت پروژه‌هایی می‌برد که تسریع نقطه‌ی عطف در آن‌ها شدنی
    است~\citep{Damodaran2012}.

  \item \emph{بازدارندگی برداشت از خزانه.} در قراردادهای مبتنی بر
    پرداخت نقطه‌ی عطف، مخارج تجمعی هر پروژه از منحنی \lr{S} خط مبنای
    هزینه پیروی می‌کند، در حالی که ورودی‌ها فقط در رویدادهای گسسته‌ی
    تأیید می‌رسند. این وضع یک گودی منفی در سرمایه‌ی در گردش می‌سازد که
    تا رسیدن دریافتی‌های نقاط عطف به مخارج ادامه دارد. هزینه‌ی بهره‌ی
    $c^r_{i,t}$ \eqref{eq:interest} این گودی را با نرخ $r/12$ در هر
    دوره به سبد تحمیل می‌کند. حتی وقتی موجودی سبد $B_t$ بزرگ باشد
    (رژیم بودجه‌ی گشاد، $\kappa \gg 1$)، عامل نمی‌تواند آزادانه بیش از
    حد تخصیص دهد، چون جلوافتادن خروجی‌های تجمعی از ورودی‌های تجمعی در
    هر پروژه هزینه‌ی نگهداری دارد که بازده تنزیل‌شده را کم می‌کند.
    به این ترتیب، انضباط سرمایه‌ی در گردش در تأمین مالی واقعی پروژه
    بدون جمله‌ی جریمه‌ی جداگانه مستقیماً در سیگنال پاداش جا می‌گیرد.
\end{enumerate}

% - - - - - - - - - - - - - - - - - - - - - - - - - - - - - -
\paragraph{بردار مشاهده}
\label{subsubsec:obs}
% - - - - - - - - - - - - - - - - - - - - - - - - - - - - - -

محیط به‌جای نمایش مستقیم حالت خام $S_t$، یک بردار مشاهده‌ی رو به
سیاست، $o_t \in \mathbb{R}^{7N+2}$، ارائه می‌دهد. این بردار از $S_t$
ساخته می‌شود، با استخراج و نرمال‌سازی سیگنال‌هایی که بیشترین ارتباط
را با تصمیم تخصیص دارند. شاخص‌های خام \lr{EVM} (یعنی $\mathrm{SPI}$،
$\mathrm{CPI}$، $\mathrm{EAC}$ و $\mathrm{TCPI}$) متغیر میانی‌اند و در
$o_t$ نمی‌آیند. در عوض، در دو گروه از ویژگی‌های عملیاتی هر پروژه
تجمیع می‌شوند.

\subparagraph{سیگنال‌های تخصیص جبرانی.} این سیگنال‌ها کمترین مخارج
لازم برای پرهیز از نقض قراردادی را رمزگذاری می‌کنند و از نگاه به آینده
روی منحنی \lr{S} محاسبه می‌شوند. در زمان پروژه‌ی $\tau_i$، کمترین
پیشرفت تجمعی لازم تا پایان دوره‌ی جاری برای ماندن در حد سقف تأخیر
چنین است:
\begin{equation}
  \underline{EV}_{i,t}
    = \max\!\bigl(0,\; PV_i(\tau_i + 1) - \delta^{\max}_i\bigr).
  \label{eq:min_prog}
\end{equation}
تخصیص جبرانی دوره‌ی جاری برابر است با:
\begin{equation}
  x^{\mathrm{cu}}_{i,t}
    = \max\!\bigl(0,\; \underline{EV}_{i,t} - EV_{i,t}\bigr)
      \cdot \mathrm{BAC}_i.
  \label{eq:catchup}
\end{equation}
نگاه به دوره‌ی بعد یک گام جلوتر می‌رود و به‌جای $PV_i(\tau_i + 1)$ از
$PV_i(\tau_i + 2)$ استفاده می‌کند:
\begin{equation}
  x^{\mathrm{cu}+1}_{i,t}
    = \max\!\bigl(0,\; PV_i(\tau_i + 2) - \delta^{\max}_i
      - EV_{i,t}\bigr) \cdot \mathrm{BAC}_i.
  \label{eq:catchup_next}
\end{equation}

\subparagraph{سیگنال‌های نقطه‌ی عطف هدف.} نقطه‌ی عطف هدف $j^*_{i,t}$
نخستین نقطه‌ی عطف تأییدنشده‌ی پروژه‌ی $i$ در دوره‌ی $t$ است. از آن چهار
سیگنال به‌دست می‌آید:
\begin{align}
  \Delta\psi^*_{i,t}
    &= \max\!\bigl(0,\; \psi_{i,j^*} - EV_{i,t}\bigr),
  \label{eq:target_pg} \\
  \Delta\tau^*_{i,t}
    &= \bigl(s_i + \tau^{\mathrm{ms}}_{i,j^*}\bigr) - t,
  \label{eq:target_tg} \\
  x^*_{i,t}
    &= \Delta\psi^*_{i,t} \cdot \mathrm{BAC}_i,
  \label{eq:target_alloc} \\
  \mathrm{NPV}^*_{i,t}
    &= \frac{\hat{P}^*_{i,t}}
            {\gamma^{\max(0,\,\Delta\tau^*_{i,t})}}
       - x^*_{i,t},
  \label{eq:target_npv}
\end{align}
که در آن $\hat{P}^*_{i,t} = \hat{P}_{i,j^*}$ پرداخت خالص نقطه‌ی عطف
هدف است. رابطه‌ی~\eqref{eq:target_npv} پرداخت مورد انتظار را به دوره‌ی
جاری تنزیل می‌کند و سرمایه‌گذاری لازم را از آن کم می‌کند. حاصل، ارزش
خالص فعلی گشودن نقطه‌ی عطف بعدی است، با این فرض که پرداخت در زودترین
دوره‌ی واجد شرایط برسد.

\subparagraph{ساخت مشاهده.} بردار مشاهده الحاق تخت‌شده و بریده‌شده‌ی
زیر است:
\begin{equation}
  o_t = \Bigl(
    \bigl\{
      \theta^{\mathrm{rem}}_{i,t},\;
      \tilde{x}^{\mathrm{cu}}_{i,t},\;
      \tilde{x}^{\mathrm{cu}+1}_{i,t},\;
      \Delta\psi^*_{i,t},\;
      \widetilde{\Delta\tau}^*_{i,t},\;
      \tilde{x}^*_{i,t},\;
      \widetilde{\mathrm{NPV}}^*_{i,t}
    \bigr\}_{i=0}^{N-1},\;
    \widetilde{\Delta B}_{t},\;
    \tilde{B}_t
  \Bigr),
  \label{eq:obs_vector}
\end{equation}
که در آن علامت تیلدا نشان‌دهنده‌ی نرمال‌سازی است: با $B_0$ برای
کمیت‌های پولی، با $T$ برای فاصله‌های زمانی، یا با
$\max_i |\mathrm{NPV}^*_{i,t}|$ برای مؤلفه‌ی \lr{NPV}. سپس مقادیر به
بازه‌های جدول~\ref{tab:obs_vector} بریده می‌شوند. برای پروژه‌های
در انتظار یا تکمیل‌شده، همه‌ی فیلدهای مربوط به پروژه صفر می‌شوند.
اسکالرهای سبد، $\widetilde{\Delta B}_t$ و $\tilde{B}_t$، به‌ترتیب
جریان نقدی خالص دوره‌ی قبل و بودجه‌ی در دسترس فعلی‌اند و هر دو با
$B_0$ نرمال شده‌اند.

\begin{table}[H]
  \centering
  \includegraphics[width=\textwidth]{tables/obs_vector.pdf}
  \caption[بردار مشاهده]{بردار مشاهده‌ی $o_t$ \eqref{eq:obs_vector}:
    مؤلفه‌ها، نرمال‌سازی و کران‌های برش. بلوک پروژه‌ی $i$ اندیس‌های
    $7i$ تا $7i+6$ را اشغال می‌کند.}
  \label{tab:obs_vector}
\end{table}

% - - - - - - - - - - - - - - - - - - - - - - - - - - - - - -
\paragraph{پیاده‌سازی سازگار با \lr{Gymnasium}}
\label{subsubsec:impl}
% - - - - - - - - - - - - - - - - - - - - - - - - - - - - - -

\lr{MDP} با نام \code{PortfolioBudgetingEnv} پیاده‌سازی شده است. این
یک محیط پیکربندی‌پذیر و سازگار با \lr{Gymnasium} است~\citep{Towers2024}
که رابط استاندارد \code{reset()} / \code{step()} را به‌همراه موارد
زیر ارائه می‌دهد:
\begin{itemize}[nosep]
  \item \code{observation\_space}: $\mathrm{Box}(7N+2)$ با کران‌های هر
    مؤلفه مطابق جدول~\ref{tab:obs_vector}؛
  \item \code{action\_space}: $\mathrm{Box}([0,1]^N)$، کسرهای خام
    تخصیص از $B_t$.
\end{itemize}

\subparagraph{دیکشنری پیکربندی.}
همه‌ی رفتار محیط با یک دیکشنری پیکربندی واحد اداره می‌شود که هنگام
مقداردهی اولیه داده می‌شود. این دیکشنری دو دسته ورودی دارد که با هم
یک آزمایش را کامل مشخص می‌کنند.

دسته‌ی اول \emph{پارامترهای ریاضی} صورت‌بندی است که در
جدول~\ref{tab:params} گردآوری شده‌اند: پارامترهای قراردادی
($\mathrm{BAC}_i$، $\pi_i$، $\phi_i$، $\psi^{\mathrm{rec}}_i$،
$\delta^{\mathrm{rec}}_i$، $\nu_i$، $w_{i,j}$)، پارامترهای زمان‌بندی
($s_i$، $D_i$، $a_i$، $b_i$، $\psi_{i,j}$،
$\tau^{\mathrm{ms}}_{i,j}$، $\delta^{\max}_i$، $\Delta^{\max}_i$،
$\theta_i$) و اسکالرهای سطح سبد ($\kappa$، $\gamma$، $r$).

دسته‌ی دوم \emph{پارامترهای ساختاری} است که در دو گروه پایینی
جدول~\ref{tab:params} آمده‌اند: اندازه‌ی سبد $N$
(\code{n\_projects})، تعداد نقاط عطف هر پروژه $M_i$
(\code{n\_milestones}) و دنباله‌ی کارایی تصادفی $\{\eta_{i,t}\}$
(\code{efficiency}) که تنها متغیر تصادفی بنیادین است. هر ورودی در هر
دو دسته به‌صورت یک خانواده‌ی توزیع با حداکثر چهار پارامتر مشخص می‌شود
و خانواده‌های \code{fixed}، \code{uniform}، \code{normal}، \code{beta}،
\code{triangular}، \code{lognormal}، \code{truncated\_normal}،
\code{categorical} و \code{custom\_ordinal} پشتیبانی می‌شوند. خانواده‌ی
\code{custom\_ordinal} یک دنباله‌ی قطعی می‌پذیرد، یعنی فهرست ساده برای
پیکربندی‌های تک‌پروژه‌ای یا فهرستی از فهرست‌ها (یکی برای هر پروژه)
برای پیکربندی‌های سبدی، و در هر اپیزود به‌ترتیب روی آن می‌چرخد. این
کار نمونه‌های کارایی را دقیقاً تکرارپذیر می‌کند و برای محک‌زنی
کنترل‌شده لازم است. برای پارامترهای اسکالر هر پروژه، فهرست ساده با
اندیس پروژه‌ی $i$ خوانده می‌شود و به هر پروژه یک مقدار ثابت از دنباله
می‌دهد. برای متغیر بنیادین کارایی $\eta_{i,t}$، اندیس به‌جای آن زمان
پروژه‌ی $\tau_i$ است و دنباله در طول دوره‌ها می‌چرخد. این طراحی اجازه
می‌دهد هر پارامتری در آزمایش‌های کنترل‌شده ثابت نگه داشته شود، یا
برای تصادفی‌سازی دامنه در هر اپیزود دوباره نمونه‌برداری
شود~\citep{Tobin2017}.

\subparagraph{چرخه‌ی عمر اپیزود.}
در هر \code{reset()}، مقدار $N$ از توزیع پیکربندی‌شده‌ی خود نمونه‌برداری
می‌شود و همه‌ی پارامترهای هر پروژه به‌طور مستقل کشیده می‌شوند. سپس
افق $T$ به‌طور خودکار از رابطه‌ی~\eqref{eq:horizon} به‌دست می‌آید.
اپیزود در این حالت‌ها تمام می‌شود: همه‌ی پروژه‌ها به وضعیت
\code{completed} یا \code{terminated} برسند، $t \geq T$ شود، یا
$B_t \leq 0$ باشد و هیچ پیش‌پرداخت در انتظاری نباشد. یک دفتر ثبت
حسابرسی \lr{SQLite} در هر گام، هر تخصیص، هر رویداد جریان نقدی، هر
شاخص \lr{EVM} و هر پرچم نقض را ثبت می‌کند، که رفتار هر سیاستی را کاملاً
قابل تفسیر و بازتولیدپذیر می‌کند.

% - - - - - - - - - - - - - - - - - - - - - - - - - - - - - -
\paragraph{سیاست‌های خط مبنای قطعی}
\label{subsubsec:baselines}
% - - - - - - - - - - - - - - - - - - - - - - - - - - - - - -

سه سیاست خط مبنای قطعی روی $o_t$ تعریف می‌شوند تا هم سطح عملکرد
پایه را مشخص کنند و هم شکاف تا کران بالای اطلاعات‌کامل را اندازه
بگیرند. هر کدام مستقیماً از سیگنال‌های جبرانی در $o_t$ می‌خوانند و
مثل هر سیاست دیگری از همان تصویرسازی سه‌مرحله‌ای
\eqref{eq:projection} می‌گذرند.

\subparagraph{رسیدن به برنامه (\lr{Reach-Plan}، $\pi^{\mathrm{rp}}$).}
دقیقاً به‌اندازه‌ای تخصیص بده که هر پروژه‌ی فعال به پیشرفت مطابق برنامه
یعنی $PV_i(\tau_i + 1)$ برسد:
\begin{equation}
  x^{\mathrm{rp}}_{i,t}
    = \max\!\bigl(0,\; PV_i(\tau_i+1) - EV_{i,t}\bigr) \cdot
      \mathrm{BAC}_i.
  \label{eq:reach_plan}
\end{equation}
این سیاست شکاف پیشرفت را در هر دوره کاملاً می‌بندد، اما زمان‌بندی نقاط
عطف، سودآوری قرارداد و کمبود بودجه را نادیده می‌گیرد.

\subparagraph{جبران \lr{T} (\lr{Catchup-T}، $\pi^{\mathrm{cu}}$).}
کمترین مقدار لازم برای پرهیز از نقض در همین دوره را تخصیص بده:
$x^{\mathrm{cu}}_{i,t}$ از رابطه‌ی~\eqref{eq:catchup}. این کم‌هزینه‌ترین
راهبردی است که از فعال‌شدن شمارنده‌ی نقض در هر دوره جلوگیری می‌کند.

\subparagraph{جبران \lr{T} بعدی (\lr{Catchup-Next-T}، $\pi^{\mathrm{cu}+1}$).}
نگاه یک‌دوره‌ای به آینده: $x^{\mathrm{cu}+1}_{i,t}$ را از
رابطه‌ی~\eqref{eq:catchup_next} تخصیص بده. این سیاست از نقض در دوره‌ی
جاری و دوره‌ی بعد هر دو جلوگیری می‌کند، به قیمت مخارج اولیه‌ی بیشتر و
بازده تنزیل‌شده‌ی کمتر.

این سه سیاست طیفی از تهاجمی (رسیدن به برنامه، بیشترین مخارج) تا
محافظه‌کارانه (جبران \lr{T}، کمترین مخارج) را می‌پوشانند و یک بازه‌ی
مبنای معنادار می‌سازند. عملکرد آن‌ها در چهار پیکربندی آزمون در
فصل~\ref{chap:evaluation} تحلیل شده است.